%%%%%%%%%%%%%%%%%%%%%%%%%%%%%%%%%%%%%%%%%%%%%%%%%%%%%%%%%%%%%%%%%%%%%%%%%%%
%%%                         Resultados                                 %%%
%%%%%%%%%%%%%%%%%%%%%%%%%%%%%%%%%%%%%%%%%%%%%%%%%%%%%%%%%%%%%%%%%%%%%%%%%%%

\chapter{Resultados esperados}
\label{ch:resultados}

Espera-se que o sistema Cognitivo atue como um agente transformador na rotina de
estudos dos alunos. O principal impacto almejado é o aumento significativo do
engajamento dos estudantes com os conteúdos de matemática voltados para o \gls{enem}.
Ao converter a resolução de questões tradicionais em um fluxo guiado pelos
pilares do \gls{pc}, o sistema pretende mitigar a resistência histórica associada à
disciplina, substituindo a postura passiva ou a frustração por uma
experiência de aprendizado ativa, fluida e psicologicamente estimulante.

A eficácia dessa abordagem lúdica será validada empiricamente por meio da
aplicação de formulários estruturados aos estudantes após os testes de uso.
Esses instrumentos de coleta de dados permitirão mensurar o impacto real da
plataforma sob duas perspectivas fundamentais, a usabilidade da interface e a
percepção do aluno diante dos problemas matemáticos. Busca-se extrair dos
questionários evidências sobre como os elementos de gamificação e o \gls{pc} podem
engajar os alunos.
